% BHU PYQ -- History: History of India (1857 - 1947) (2025-26 NEP)
\documentclass[11pt,a4paper]{article}
\usepackage[a4paper,top=2.5cm,bottom=2.5cm,left=2.8cm,right=2.8cm,headheight=14pt]{geometry}
\usepackage{amsmath,amssymb}
\usepackage{fontspec}
\setmainfont{Kohinoor Devanagari}
\usepackage{microtype,fancyhdr,enumitem,hyperref,lastpage,parskip}

\hypersetup{colorlinks=true,urlcolor=black,linkcolor=black,pdftitle={HISMJ32: History of India (1857 - 1947) -- BHU NEP 2025-26}}

\pagestyle{fancy}\fancyhf{}
\renewcommand{\headrulewidth}{0.4pt}\renewcommand{\footrulewidth}{0.4pt}
\fancyhead[L]{\small \textit{Banaras Hindu University}}
\fancyhead[R]{\small \textit{HISMJ32: History of India (1857 - 1947) (2025-26)}}
\fancyfoot[C]{\small Page \thepage\ of \pageref{LastPage}}

\newlist{parts}{enumerate}{2}
\setlist[parts,1]{label=(\alph*),leftmargin=*,itemsep=4pt,topsep=3pt}
\setlist[parts,2]{label=(\roman*),leftmargin=*,itemsep=2pt,topsep=2pt}

\newcommand{\pts}[1]{\hfill \textit{[#1]}}

\begin{document}

\begin{center}
  {\large\bfseries BANARAS HINDU UNIVERSITY}\\[2pt]
  {\normalsize B. A. (Hons.) Semester III Examination, 2025-26 (NEP)}\\[6pt]
  \rule{0.8\linewidth}{0.5pt}\\[6pt]
  {\large\bfseries HISTORY}\\[4pt]
  {\large\bfseries Paper Code: HISMJ32 : History of India (1857 -- 1947)}\\[4pt]
  {\normalsize Time: 2:30 Hours \hfill Full Marks: 70}
\end{center}
\rule{\linewidth}{0.4pt}
\smallskip

\noindent \textit{Instructions: This question paper comprises three Sections A, B and C. All questions are compulsory. Write your Roll No.\ at the top immediately on receipt of this question paper.}

\smallskip
\noindent \textit{इस प्रश्न-पत्र में तीन खण्ड अ, ब तथा स हैं। सभी प्रश्न अनिवार्य हैं।}

\bigskip

\section*{SECTION -- A / खण्ड -- अ}
\noindent \textit{Note: Objective Type Questions. Write the answer of each question in not more than 50 words.}\pts{5 $\times$ 2 = 10}

\noindent \textit{वस्तुनिष्ठ प्रश्न। प्रत्येक प्रश्न का उत्तर 50 शब्दों से अधिक न हो।}

\begin{enumerate}[label=\textbf{\arabic*.},leftmargin=*,itemsep=8pt]
  \item 
  \begin{parts}
    \item Write the provisions of the Indian Councils Act of 1892.\pts{2}
    
    भारतीय परिषद अधिनियम, 1892 के प्रावधान लिखिए।

    \item What was the Montagu-Chelmsford Reforms (1919)?\pts{2}
    
    मांटेग्यू-चेम्सफोर्ड सुधार क्या था?

    \item Write a brief note on the Arya Samaj.\pts{2}
    
    आर्य समाज पर एक संक्षिप्त टिप्पणी लिखिए।

    \item What is meant by the Revolutionary Movement in India?\pts{2}
    
    भारत में क्रांतिकारी आंदोलन से क्या अभिप्राय है?

    \item What was the Hindu Widow Remarriage Act of 1856?\pts{2}
    
    हिंदू विधवा पुनर्विवाह अधिनियम, 1856 क्या था?
  \end{parts}
\end{enumerate}

\bigskip

\section*{SECTION -- B / खण्ड -- ब}
\noindent \textit{Note: Short Answer Type Questions. Answer each in about 250 words.}\pts{3 $\times$ 10 = 30}

\noindent \textit{लघु उत्तरीय प्रश्न। प्रत्येक प्रश्न का उत्तर लगभग 250 शब्दों में दीजिए।}

\begin{enumerate}[label=\textbf{\arabic*.},leftmargin=*,start=2,itemsep=12pt]

  \item Critically examine the significance of Queen Victoria's Proclamation of 1858.\pts{10}
  
  1858 के महारानी विक्टोरिया के घोषणा-पत्र के महत्त्व का समालोचनात्मक परीक्षण कीजिए।

  \begin{center}\textbf{OR / अथवा}\end{center}

  Critically examine Ishwar Chandra Vidyasagar's role in the campaign for Widow Remarriage in mid-19th-century Bengal.\pts{10}
  
  19वीं शताब्दी के मध्य बंगाल में विधवा पुनर्विवाह आंदोलन में ईश्वर चंद्र विद्यासागर की भूमिका का समालोचनात्मक परीक्षण कीजिए।

  \item Explain the formation and activities of the Congress Socialist Party (CSP).\pts{10}
  
  कांग्रेस सोशलिस्ट पार्टी के गठन तथा गतिविधियों की व्याख्या कीजिए।

  \begin{center}\textbf{OR / अथवा}\end{center}

  Write an essay on Swami Dayanand Saraswati and the Arya Samaj.\pts{10}
  
  स्वामी दयानंद सरस्वती एवं आर्य समाज पर एक निबंध लिखिए।

  \item Discuss the Non-Cooperation Movement launched by Mahatma Gandhi and its impact on the Indian Freedom Struggle.\pts{10}
  
  महात्मा गांधी द्वारा चलाए गए असहयोग आंदोलन तथा भारतीय स्वतंत्रता संग्राम पर इसके प्रभाव की विवेचना कीजिए।

  \begin{center}\textbf{OR / अथवा}\end{center}

  Discuss the causes and significance of the Quit India Movement of 1942.\pts{10}
  
  1942 के भारत छोड़ो आंदोलन के कारणों एवं महत्त्व की विवेचना कीजिए।

\end{enumerate}

\bigskip

\section*{SECTION -- C / खण्ड -- स}
\noindent \textit{Note: Long Answer Type Questions. Answer each in about 500 words.}\pts{2 $\times$ 15 = 30}

\noindent \textit{दीर्घ उत्तरीय प्रश्न। प्रत्येक प्रश्न का उत्तर लगभग 500 शब्दों में दीजिए।}

\begin{enumerate}[label=\textbf{\arabic*.},leftmargin=*,start=5,itemsep=14pt]

  \item ``The first half of the twentieth century witnessed the transformation of Indian nationalism into a mass political movement.'' Discuss.\pts{15}
  
  ``बीसवीं शताब्दी के पूर्वार्द्ध में भारतीय राष्ट्रवाद का जन-राजनीतिक आंदोलन के रूप में रूपांतरण हुआ।'' विवेचना कीजिए।

  \begin{center}\textbf{OR / अथवा}\end{center}

  Critically examine the development of revolutionary activities in India and abroad from the late 19th century to 1947.\pts{15}
  
  19वीं शताब्दी के उत्तरार्द्ध से 1947 तक भारत एवं विदेशों में क्रांतिकारी गतिविधियों के विकास का समालोचनात्मक परीक्षण कीजिए।

  \item Trace the events leading to the Partition and Independence of India in 1947.\pts{15}
  
  1857 से 1947 के मध्य भारत के विभाजन तथा स्वतंत्रता की ओर ले जाने वाली प्रमुख घटनाओं का विवरण दीजिए।

  \begin{center}\textbf{OR / अथवा}\end{center}

  Discuss the role of Mahatma Gandhi in shaping the strategy and mass base of the Indian National Movement.\pts{15}
  
  भारतीय राष्ट्रीय आंदोलन की रणनीति एवं जन-आधार को आकार देने में महात्मा गांधी की भूमिका की विवेचना कीजिए।

\end{enumerate}

\end{document}
